% ============================================================
% ch10：讨论与局限（对应论文 §7 Discussion and limitations）
% ============================================================

\section{讨论与局限}\label{sec:ch10}

\begin{keybox}{本章速览}
\begin{itemize}
  \item 可交付的监测器：$\lhatg=\mathbf{g}^{\!\top}H\mathbf{g}/
        \mathbf{g}^{\!\top}\mathbf{g}$（单次 HVP）在 EoS 对齐区
        （$|\cos(\mathbf{g},\vtope)|>0.7$）近似精确；绝对值带
        ${\sim}17\%$ 瑞利下界松弛（相对误差 $0.169\pm0.026$），更高精度
        用热启动幂迭代（${\sim}5$ HVPs）；低于对齐区用 $\|\mathbf{g}\|$
        兜底。
  \item 控制的负结果本身就是建议：别指望秩-1/proxy 驱动的曲率控制器把
        $\lmax$ 压到 EoS 均衡之下——predictor$\ne$actuator 是结构性的
        （Claim 1 / Proposition 1），不是调参失败；要压只有全局动作
        （SAM）。
  \item 失败是 sharpness 特异的、不是优化全面的：精确 $\vtope$
        的 rank-1/rank-$k$ 预条件器最终 loss 更低（$0.090$ vs
        $0.111$），但去的是更尖的最小值。
  \item 七条局限逐条在案，其中第 (vi) 条是反直觉亮点：上逃逸臂
        （$\lmax=261.8\pm5.4$）的 held-out 测试 MSE $0.0932$，竟优于
        钉住的 SGD（$0.1142$）与 SAM（$0.1296$）——最尖的臂泛化最好。
\end{itemize}
\end{keybox}

\textbf{本章在论文中的位置。}论文 \S7 是“讨论与局限”：先给实践建议
（监测器）、再划清“它不适合什么”、最后逐条列出七条局限 (i)--(vii)。这
个顺序本身就是论文的论点结构：\emph{能交付的先交付，交付不了的讲清楚为
什么交付不了，最后诚实地圈出适用域}。本章按同样顺序解读，并特别指出：
这份局限清单不是防守姿态，而是一份方法论示范——几乎每条局限都对应一个
已做的补救实验或一个明确的开放性声明。

\subsection{可交付的监测器：什么时候信 $\lhatg$、信到什么程度}\label{sec:ch10-monitor}

论文的实践建议浓缩为一句话：$\lhatg$ 是实用的 $\lmax$ 监测器，但要按区
域校准信任度。

\begin{itemize}
  \item \textbf{EoS 对齐区（核心使用场景）}：当
        $|\cos(\mathbf{g},\vtope)|>0.7$ 时，$\lhatg$ 与 $\lmax$
        的去趋势相关近乎精确；而渐进尖锐化本身会把对齐推到
        $0.85$--$0.9$——也就是说，\emph{你越是在 EoS 体制里，这个单 HVP
        监测器越准}。成本：每次测量一次 Hessian-向量积。
  \item \textbf{绝对值的精度}：作为沿梯度方向的瑞利商，$\lhatg$ 天然是
        $\lmax$ 的下界，带 ${\sim}17\%$ 的松弛（相对误差
        $0.169\pm0.026$）。要压到 $17\%$ 以内，用热启动幂迭代
        （${\sim}5$ HVPs；Ramasinghe 等人，2026）。
  \item \textbf{低于 EoS 对齐时（早训、归一化主导体制）}：$\lhatg$ 会
        低估，此时原始 $\|\mathbf{g}\|$ 是粗但架构无关的兜底指标。
\end{itemize}

\textbf{控制端的建议是否定式的}：不要期待秩-1 或 proxy 驱动的曲率控制
器把 $\lmax$ 压到 EoS 均衡之下——第~\ref{sec:ch09}~章的命题 1 说明这是
结构性的不可能（对缩放型单方向动作），不是调参问题。如果一定要压尖锐
度，只有全局动作（SAM 一族）可用。这些建议共同构成论文标题的落地形式：
“预测器，不是执行器”（predictor, not actuator）。

\subsection{它不适合什么：结构性失败与它的精确边界}\label{sec:ch10-notfor}

\textbf{predictor$\ne$actuator 之分是结构性的。}论文反复强调：这不是
“没调好”——Claim 1 与命题 1（第~\ref{sec:ch09}~章）说明，阻尼顶部方向
把帽向上移（上逃逸）、阻尼正交补把帽留在原地、梯度方向传感器因不精确
（而非滞后）钉在均衡上。三种单方向干预、三种失败方式，全部可由同一
二维递推预言。只有全局动作向下。

\textbf{失败是 sharpness 特异的，不是优化全面的。}这是一个容易漏读、但
对优化器研究者很关键的限定：精确 $\vtope$ 的 rank-1/rank-$k$
预条件器最终 loss \emph{更低}（$0.090$ vs SGD 的 $0.111$）——它是个不
坏的优化器！只是它抵达的是一个\emph{更尖}的最小值（上逃逸值
$261.8\pm5.4$），与“压尖锐度”的控制目标背道而驰。换言之：\emph{逃出
EoS 不是好优化的必要条件，只是尖锐度控制本身的必要条件}。把这两个目标
混为一谈，就会误把上逃逸臂读成“更差的优化器”。

\textbf{对优化器设计空间的净效应是“收窄”}：曲率控制路线在秩 1 处关闭
（单方向缩放型动作全部失败），剩余可行的是密集的步长调谐线——CDAT
（Roulet 等人，2024）、warmup 调度器（Kalra 等人，2024）——它们用更新
方向调步长，不做谱执行。结合第~\ref{sec:ch08}~章的实证（SAM 与 rank-1
朝相反方向移动阈值），“全局动作 vs 步长调谐”就是论文留给优化器设计
的两条活路。

\subsection{七条局限逐条解读}\label{sec:ch10-limits}

\begin{itemize}
  \item \textbf{(i) 规模与架构。}控制实验（C1、C3、rank-$k$、X-SAM、
        多尖锐）确实源自 MLP；但核心三向控制（SGD vs rank-1 plant vs
        SAM）与 rank-$k$ 扫描已在标准 ResNet-18（11.2M）/CIFAR-100 上
        复现：双结局分裂（第~\ref{sec:ch08}~章）；C2（更小 ResNet 上的
        $\gamma$-cap）是第二个非 MLP 控制。ImageNet/LLM 规模验证开放。
        特别值得学习的是论文对\emph{两个测量警告}的处理：其一，早期
        ResNet-18 跑用了默认（非确定性）CUDA，论文测得 seed 内非确定性
        会产生 ${\sim}6$ 单位的伪 $\lmax$ 差距（确定性模式消除之），
        因此头条三向改用确定性 CUDA + ARPACK（$n=5$）重报：SAM
        $22.63\pm1.73$ vs SGD/rank-1 ${\sim}29$、CI 不重叠；其二，原
        rank-1 plant 的惰性追溯到幂迭代感知伪收敛（对真实
        $\lmax\approx36$ 的 $+45\%$ 低估），不是体制属性，于是论文造了
        true-$\lmax$ plant（ARPACK 感知、无钳位下限）以真实力开火（除数
        中位 $1.86$）——尖锐度轨迹仍与 SGD 重叠（$n=5$，有动作仍被
        钉住）。两个警告都\emph{被正面化解}，而不是被埋进脚注。
  \item \textbf{(ii) 形式化程度。}核心代数主张形式化为命题 1：无翻号
        周期-2 轨道 $y^{*}\le0$、条件 $y$-收缩
        $|\mu_y|<1\iff(1-\gamma)\eta^2\alpha<1$。这是\emph{二维约化
        定理}，不是完整轨迹定理：多特征值解耦（G1）、旋转奇点（G2）、
        时变 $\gamma$（G3）开放；且控制论框架（观测器-控制器分离原理）
        是从线性控制借来的\emph{类比}，而这里的动力学是非线性、随机的。
  \item \textbf{(iii) 只测了 vanilla SAM。}ASAM、SAF 等变体共享全局
        动作结构，但本文的负结果（rank-1 无效）\emph{不外推}到它们——
        对它们的任何判断都需要单独实验。
  \item \textbf{(iv) rank-$k$ 的证据结构。}$k$-不变出现在一维尖锐默认
        几何与 ResNet-18（经钳位休眠）上；$k$-依赖出现在多尖锐构造上
        （每个被解除帽的方向独立上逃逸）；但\emph{任何}测试过的 $k$、
        在\emph{任何}景观上都没有下逃逸。密尖景观 + 足够大 $k$ 能否下
        逃逸，开放。
  \item \textbf{(v) 帽机制的体制适用域。}自稳定帽机制在全批 GD 体制
        精确；动量/小批的 ResNet-18 结果对照重球/随机 EoS 阈值
        （Andreyev 等人，2024）做\emph{定性}处理——排序与 CI 分离成立，
        精确钉住不主张（第~\ref{sec:ch08}~章的双结局读法）。
  \item \textbf{(vi) 上逃逸最小值不付测试代价（反直觉亮点）。}评估
        上逃逸臂（$k=1$，$\lmax^{\mathrm{final}}=261.8\pm5.4$）在同一
        teacher 的 held-out 数据上（$n_{\mathrm{test}}=8192$、$n=5$
        种子、贝叶斯最优地板 $0.09$——训练数据生成时注入了 std $0.3$
        的标签噪声，其方差 $0.3^{2}=0.09$ 是任何模型都不可能突破的
        测试 MSE 下限：“贝叶斯最优”指即使完全知道数据生成真相的最优
        预测者也只能达到这个水平）：测试 MSE $0.0932$，vs 钉在
        $2/\eta=40$ 的 SGD 的 $0.1142$、vs SAM（$\lmax\approx21$）的
        $0.1296$——\emph{最尖的臂泛化最好}，方向在所有 seed 上一致
        （这三个 MSE 论文未附 CI，但``5 个种子全部同向''本身已是较强
        旁证：若两臂真无差异，全部同向的概率约
        $2\times(1/2)^{5}=1/16\approx6\%$），
        逆着朴素的平坦度-泛化预测（Keskar 等人，2017）。且逃逸臂的
        训练 MSE（$0.090$）已坐在噪声地板上，所以优势来自更深的信号
        拟合、不是测试侧过拟合——地板 $0.09$ 是纯噪声方差：训练 MSE
        若\emph{低于} $0.09$ 才说明模型在死记噪声，恰好停在 $0.09$
        即信号全拟合、噪声零拟合。论文明确把这一发现的适用域限定在合成
        teacher--student 任务——它不推翻平坦度-泛化叙事，但提醒：
        “更尖”不自动等于“更差”。
  \item \textbf{(vii) 训练预算与外部效度。}规模实验用 5000 步
        （${\approx}13$ epochs）、恒定 lr、无增强——一个欠拟合体制
        （CIFAR-10 ${\sim}80\%$、CIFAR-100 ${\sim}52\%$，远低于标准
        200-epoch 增强配方的峰值）。这是\emph{有意取舍}：确定性 CUDA、
        密集 $\lmax$ 测量、$n\ge5$（X-SAM $n=10$）种子，买种子级可复现
        与测量密度，付峰值精度。所以机制主张是“\emph{同预算下优化器间
        的排序}”，不是精度主张。首个更长日程复现——12000 步
        （${\approx}31$ epochs）+ pad/crop + 翻转、$n=2$ 种子/臂——
        保序：SAM 晚期 $\lmax$ $18.1$ vs SGD $21.9$ / rank-1 $22.8$，
        SAM 精度 $0.68$ vs $0.66$（第~\ref{sec:ch08}~章）；完整
        200-epoch 配方未测。所有对比报告 95\% CI，但无正式功效分析（power analysis，事先计算检出效应所需样本量的统计设计）；
        功效不足的配置（$n\le3$）被标记为仅作支持性观察。
\end{itemize}

\subsection{解读：这份局限清单本身是方法论示范}\label{sec:ch10-reading}

把七条放在一起看，会发现它们不是同质的“免责声明”，而分三类：

\begin{itemize}
  \item \textbf{已被补救的}：(i) 的两个测量警告各自催生了一个更强的
        复现（确定性协议、true-$\lmax$ plant），(vii) 的欠拟合质疑催生
        了 12000 步增强复现——\emph{质疑进来，实验顶上}。
  \item \textbf{已划适用域的}：(iii)(v)(vi) 不辩护、只圈地——负结果
        不外推到未测变体；机制只在 GD 体制精确；上逃逸泛化优势只在
        合成任务上主张。
  \item \textbf{明确开放的}：(ii)(iv) 与 G1--G5——理论缺口与实验缺口
        都点名，成为后续工作的问题清单。
\end{itemize}

对一个以负结果为主线的论文来说，这个结构是可信度的来源：读者可以逐条
检查“这条局限有没有对应的动作”，而论文对每条都有交代。这与
第~\ref{sec:ch08}~章的伪收敛自纠、第~\ref{sec:ch09}~章的“命题不主张
什么”注记，共同构成这篇论文的方法论签名。

\subsection*{本章要点}
\begin{itemize}
  \item 监测器是可交付的：$\lhatg$ 单 HVP、EoS 对齐区
        （$|\cos|>0.7$，实际达 $0.85$--$0.9$）近似精确、绝对值
        ${\sim}17\%$ 松弛（$0.169\pm0.026$）；更高精度 ${\sim}5$ HVPs
        热启动幂迭代；对齐区之外用 $\|\mathbf{g}\|$ 兜底。
  \item 控制建议是否定式且结构性的：单方向缩放型动作（秩-1/正交/滞后
        传感器）全部失败并可由二维递推预言；要压 $\lmax$ 只有全局动作
        （SAM）；剩余可行路线是步长调谐（CDAT、warmup）。
  \item 失败是 sharpness 特异的：rank-1/rank-$k$ 预条件器最终 loss
        $0.090<0.111$（SGD）但更尖——“逃出 EoS”与“好优化”是两个
        独立目标。
  \item 局限 (vi) 是亮点：最尖的臂（$261.8\pm5.4$）测试 MSE 最好
        （$0.0932$ vs $0.1142$ vs $0.1296$，地板 $0.09$，方向跨 seed
        一致），限定于合成 teacher--student。
  \item 局限清单分三类（已补救/已圈地/明确开放），构成论文的方法论
        示范——每条局限都有对应的实验动作或开放性声明。
\end{itemize}
